% Table tab:viz_folds_methods — Conclusive / Inconclusive verdict bands
% over the five-fold patient-stratified CV protocol.
% Wide table (4 cols, last col long): use tabularx so the "Genes" column
% wraps cleanly inside the text width.
\small
\setlength{\tabcolsep}{6pt}
\begin{tabularx}{\linewidth}{l c c X}
\toprule
\textbf{Verdict} & \textbf{Mean AUROC band} &
  \makecell{\textbf{Per-slide score}\\\textbf{still reported?}} &
  \textbf{Genes (this cohort)} \\
\midrule
Conclusive   & $\geq 0.70$ & Yes & \emph{none} \\
Inconclusive & $< 0.70$    & Yes &
  \textit{TP53} ($0.661$), \textit{STK11} ($0.608$),
  \textit{EGFR} ($0.582$), \textit{RBM10} ($0.552$),
  \textit{KRAS} ($0.517$), \textit{KEAP1} ($0.488$) \\
\bottomrule
\end{tabularx}
\begin{flushleft}\footnotesize
The verdict is a population-level signal-presence flag, not a
per-slide confidence threshold; the per-slide model probability is
reported independently in every case.
Values are the best of the four conditions per gene under
patient-preserved folds, pooled out-of-fold and averaged over five seeds
(\cref{tab:mutation_luad_resultssec}).
\par\smallskip
Under nested cross-validation on the corrected $458$-slide cohort,
\textbf{no LUAD target clears the threshold}: the highest value in the
study is \textit{TP53} at $0.661$, and the band is empty. Earlier
single-split estimates on the superseded inventory placed \textit{TP53}
and \textit{EGFR} above $0.70$; removing early stopping on the scored
fold and excluding squamous and solid-tissue-normal sections removes that
margin. The Conclusive band is populated in this work only by the other
two modalities---glioma IDH and 1p/19q, and both electrocardiogram
targets (\cref{tab:crossdomain_merged})---which is the intended
behaviour of a verdict that is allowed to come back empty.
\end{flushleft}
